%%%%%%%%%%  scan idx 494  |  folio 481  %%%%%%%%%%
\SGAPrintPage{481}%
Applying
$i^*\circ j'^*$ gives the exact sequence:
\[
0\longrightarrow (p''\circ j'\circ i)^*(\underline{\Omega}^1_{V'/Z})\longrightarrow (j'\circ i)^*(\underline{\Omega}^1_{V''/Z})\longrightarrow (j'\circ i)^*(\underline{\Omega}^1_{V''/V'})\longrightarrow0.
\]
Since $j'^*(\underline{\Omega}^1_{V''/V'})=\underline{\Omega}^1_{V/Y}$ by hypothesis (G), we obtain:
\sgaStableTarget{sga6.e08.eq.2-6-2}\begin{equation}\tag{2.6.2}
0\longrightarrow f^*(j^*(\underline{\Omega}^1_{V'/Z}))\longrightarrow (j'\circ i)^*(\underline{\Omega}^1_{V''/Z})\longrightarrow i^*(\underline{\Omega}^1_{V/Y})\longrightarrow0.
\end{equation}

Combining (2.6.1) and (2.6.2), we are reduced to checking that the diagram:
\[
\begin{tikzcd}[column sep=large,row sep=large]
0 \arrow[r] & f^*(\underline{J}/\underline{J}^2) \arrow[r] \arrow[d,"f^*(d)"'] & \underline{K}/\underline{K}^2 \arrow[r] \arrow[d,"d"'] & \underline{I}/\underline{I}^2 \arrow[r] \arrow[d,"d"'] & 0\\
0 \arrow[r] & f^*(j^*(\underline{\Omega}^1_{V'/Z})) \arrow[r] & (j'\circ i)^*(\underline{\Omega}^1_{V''/Z}) \arrow[r] & i^*(\underline{\Omega}^1_{V/Y}) \arrow[r] & 0
\end{tikzcd}
\]
is commutative, which is done locally without difficulty by returning to the definitions of the morphisms involved.

From the resulting exact sequence of complexes, one deduces the announced distinguished triangle in the derived category.

\sgaStableTarget{sga6.e08.cor.2-7}\textbf{Corollary 2.7.} \uline{With the notation and hypotheses of} 2.6, \uline{we have in} $K^{\bullet}(X)$ \uline{the following equality}:
\[
T_{g\circ f}=T_f+f^*(T_g).
\]

\sgaStableTarget{sga6.e08.sec.3}\subsection*{3. Riemann--Roch theorem: statement}

\sgaStableTarget{sga6.e08.para.3-1}\textbf{3.1.} In section 3, $Y$ is a quasi-compact scheme having an ample sheaf.

Let $f:X\to Y$ be a projective complete-intersection morphism (1.1); suppose in addition that the virtual relative dimension of $f$ (1.9)

%%%%%%%%%%  scan idx 495  |  folio 482  %%%%%%%%%%
\SGAPrintPage{482}%
is constant on $X$. Since $f$ is projective, there exists a closed immersion $X\to P$, where $P$ is a projective bundle over $Y$, and since $Y$ has an ample sheaf, one may suppose that $P=\mathbf{P}^r_Y$. We thus have a factorization:
\[
\begin{tikzcd}[column sep=large,row sep=large]
X \arrow[rr,"i"] \arrow[dr,"f"'] && \mathbf{P}^r_Y \arrow[dl,"g"]\\
& Y. &
\end{tikzcd}
\]
Since $f$ is a complete intersection, the immersion $i$ is regular (1.2).

By 1.7, $f$ is a perfect morphism. Moreover, $Y$ has an ample sheaf, and consequently so do $X$ and $P$. One can then define an additive homomorphism:
\[
f_*:K^{\bullet}(X)\longrightarrow K^{\bullet}(Y)
\]
(IV 2.12), the $K^{\bullet}$ defined in terms of locally free sheaves of finite type, or of perfect complexes, being identical. One can also define additive homomorphisms:
\[
i_*:K^{\bullet}(X)\longrightarrow K^{\bullet}(P);
\qquad
 g_*:K^{\bullet}(P)\longrightarrow K^{\bullet}(Y),
\]
and we have: $f_*=g_*\circ i_*$.

Let $d$ be the virtual relative dimension of $X$ over $Y$. The codimension of $i$ is then $r-d$. By VII 4.6, the homomorphism $i_*$ gives, after tensoring with $\mathbb Q$, inclusions:
\[
i_*\bigl(\Fil^kK^{\bullet}(X)_{\mathbb Q}\bigr)
\subset
\Fil^{k+r-d}K^{\bullet}(P)_{\mathbb Q}.
\]
By VI 5.8, the homomorphism $g_*$ gives inclusions:
\[
g_*\bigl(\Fil^{k'}K^{\bullet}(P)\bigr)
\subset
\Fil^{k'-r}K^{\bullet}(Y).
\]
Composing the two, we thus obtain:

%%%%%%%%%%  scan idx 496  |  folio 483  %%%%%%%%%%
\SGAPrintPage{483}%
\sgaStableTarget{sga6.e08.prop.3-2}\textbf{Proposition 3.2.} \uline{Let} $f:X\to Y$ \uline{be a projective morphism}, \uline{a complete intersection}, \uline{whose virtual relative dimension} $d$ \uline{is constant}. \uline{Suppose that} $Y$ \uline{is quasi-compact}, \uline{and has an ample sheaf}. \uline{Then for every} $k\in\mathbb Z$ \uline{we have}:
\[
f_*\bigl(\Fil^kK^{\bullet}(X)_{\mathbb Q}\bigr)
\subset
\Fil^{k-d}K^{\bullet}(Y)_{\mathbb Q}.
\]

We thus obtain, in particular, an additive homomorphism:
\[
f_{\mathrm{gr}}:\Gr K^{\bullet}(X)_{\mathbb Q}\longrightarrow
\Gr K^{\bullet}(Y)_{\mathbb Q}
\]
sending $\Gr^kK^{\bullet}(X)_{\mathbb Q}$ into $\Gr^{k-d}K^{\bullet}(Y)_{\mathbb Q}$.

\sgaStableTarget{sga6.e08.para.3-3}\textbf{3.3.} Recall that in (V 6.1) we defined a functor $A\longmapsto \operatorname{Ch}(A)$ on the category of graded $K$-algebras, with values in the category of augmented $K$-$\lambda$-algebras, where $K$ is a binomial ring (V 2.7). If $A$ is a graded $K$-algebra, we have:
\[
\operatorname{Ch}(A)=K\times(1+\widehat A^+).
\]

For every augmented $K$-$\lambda$-algebra $\Lambda$, we defined a $K$-$\lambda$-homomorphism, called the ``completed Chern class'' (V 6.7):
\[
\widetilde c=(\varepsilon,c):\Lambda\longrightarrow\operatorname{Ch}(\Gr_{\lambda}\Lambda).
\]

Finally, we defined a ring homomorphism, functorial in $A$, called the Chern character (V 6.4):
\[
\ch:\operatorname{Ch}(A)\longrightarrow K\oplus\widehat A^+_{\mathbb Q},
\]
and a group endomorphism, called the Todd character (V 1.16):
\[
\Todd:1+\widehat A^+\longrightarrow1+\widehat A^+.
\]

%%%%%%%%%%  scan idx 497  |  folio 484  %%%%%%%%%%
\SGAPrintPage{484}%
We shall apply these definitions to the augmented $\lambda$-rings $K^{\bullet}(X)$ and $K^{\bullet}(Y)$, the augmentation rings $H^0(X,\underline{\mathbb Z})$ and $H^0(Y,\underline{\mathbb Z})$ being binomial (V 2.8). To simplify notation, we shall put, for every scheme $Z$ and every $z\in K^{\bullet}(Z)$:
\[
\ch(z)=\ch(\widetilde c(z));
\qquad
\Todd(z)=\Todd(c(z)).
\]
For every graded algebra $A$, we shall regard the group $1+\widehat A^+$ as contained in $\widehat A$. Finally, given a homomorphism $f:A\to B$ which is additive and compatible with the gradings up to a shift, we shall still denote by $f$ the homomorphism $\widehat A\to\widehat B$ obtained from it.

\sgaStableTarget{sga6.e08.para.3-4}\textbf{3.4.} Let $X$ be a scheme. For every locally free $\underline{O}_X$-Module of finite type $\underline{E}$, denote by $\check E$ the class in $K^{\bullet}(X)$ of the dual $\underline{\check E}$ of $\underline{E}$. It is clear that the function associating $\check E$ to $\underline{E}$ is additive on the category of locally free $\underline{O}_X$-Modules of finite type, and therefore defines an endomorphism of $K^{\bullet}(X)$, which we shall denote by $x\rightsquigarrow\check x$.

\sgaStableTarget{sga6.e08.prop.3-5}\textbf{Proposition 3.5.} \uline{Let} $X$ \uline{be a quasi-compact scheme}. \uline{Then for every} $x\in K^{\bullet}(X)$ \uline{and every} $i\geq0$, \uline{we have}:
\[
c^i(\check x)=(-1)^i c^i(x).
\]

The relation to be proved may also be written:
\[
c_t(\check x)=c_{-t}(x),
\]
where we put: $c_t(x)=\sum_{i\geq0}c^i(x)t^i$.

By additivity, one may therefore suppose that $x$ is the class of a locally free $\underline{O}_X$-Module of finite type $\underline{E}$. Moreover, after replacing $X$ by the flag scheme of $\underline{E}$, which gives an injection on the

%%%%%%%%%%  scan idx 498  |  folio 485  %%%%%%%%%%
\SGAPrintPage{485}%
$\Gr K^{\bullet}$ (VI 5.5), one
may suppose that $\underline{E}$ is an invertible sheaf. We then have: $\underline{\check E}=\underline{E}^{-1}$; hence $\check x=x^{-1}$.

The $c^i(\check x)$, $c^j(x)$ vanish for $i,j\geq2$. It therefore suffices to show that $c^1(x^{-1})=-c^1(x)$. Now $c^1(x)$ is the class in $\Gr^1K^{\bullet}(X)$ of $x-1$, and $c^1(x^{-1})$ that of $x^{-1}-1$. The equality:
\[
x^{-1}-1=1-x+(x-1)(1-x^{-1})
\]
gives:
\[
x^{-1}-1\equiv1-x\qquad\mathrm{mod.}\ \Fil^2K^{\bullet}(X);
\]
the result follows.

\sgaStableTarget{sga6.e08.thm.3-6}\textbf{Theorem 3.6. (Riemann--Roch).} \uline{Let} $f:X\to Y$ \uline{be a projective complete-intersection morphism} (1.1), \uline{whose virtual relative dimension} $d$ \uline{is constant} (1.9). \uline{Suppose that} $Y$ \uline{is quasi-compact}, \uline{and has an ample sheaf}. \uline{With the notation of} 3.3, \uline{we then have for every} $x\in K^{\bullet}(X)$ \uline{the relation}:
\[
\ch(f_*(x))=f_{\mathrm{gr}}\bigl(\ch(x)\cdot\Todd(\check T_f)\bigr),
\]
$\check T_f$ \uline{denoting the dual element of the class of the relative cotangent complex} $\underline{L}^{X/Y}_\bullet$ (2.5 \uline{and} 2.3).

In other words, the diagram
\[
\begin{tikzcd}[column sep=huge,row sep=large]
K^{\bullet}(X) \arrow[r,"\ch\circ\widetilde c"] \arrow[d,"f_*"'] & \Gr K^{\bullet}(X)_{\mathbb Q} \arrow[d,"x\mapsto f_{\mathrm{gr}}(x\cdot\Todd(\check T_f))"]\\
K^{\bullet}(Y) \arrow[r,"\ch\circ\widetilde c"] & \Gr K^{\bullet}(Y)_{\mathbb Q}
\end{tikzcd}
\]
is commutative.
